\chapter{Error Handling}

\section{Error model}

\hayashi{} distinguishes two types of error:

\begin{description}
  \item[Parse errors] Detected before execution, while analysing the code.
        Include line number and code snippet.
  \item[Runtime errors] Occur during execution — undefined variable, wrong
        type, division by zero, missing key.
\end{description}

Both types display the code snippet and line number, making the source easy
to locate:

\begin{lstlisting}[style=inline]
error: line 5: undefined variable 'income' — did you mean 'gross_income'?
   5 | display income * 12
     | ^^^^^^^^^^^^^^^^^^^
\end{lstlisting}

\section{\texttt{try} / \texttt{catch}}

Runtime errors can be caught with \lstinline{try}/\lstinline{catch}:

\begin{lstlisting}[caption={Basic try/catch}]
try {
  let result = 10 / 0
  display result
} catch e {
  display f"Caught error: {e}"
}
\end{lstlisting}

The variable \lstinline{e} inside \lstinline{catch} contains the error
message as a string.

\subsection*{Catching file errors}

\begin{lstlisting}
try {
  load "missing_file.csv" as df
} catch e {
  display f"File not found: {e}"
}
\end{lstlisting}

\subsection*{\texttt{try} as an expression}

\lstinline{try}/\lstinline{catch} can return a value:

\begin{lstlisting}
let value = try {
  int("not a number")
} catch _ {
  -1
}
display value    // -1
\end{lstlisting}

\section{Common errors and their messages}

\begin{center}
\begin{tabular}{ll}
\toprule
\textbf{Situation} & \textbf{Typical message} \\
\midrule
Undeclared variable   & \texttt{undefined variable 'x' — did you mean 'y'?} \\
Unknown function      & \texttt{undefined function 'foo' — did you mean 'for'?} \\
Wrong type            & \texttt{type mismatch — expected Int, got String} \\
Index out of bounds   & \texttt{index 5 out of bounds (len=3)} \\
Missing dict key      & \texttt{key 'name' not found in dict} \\
Integer division by zero & \texttt{division by zero} \\
File not found        & \texttt{file not found: 'path.csv'} \\
\bottomrule
\end{tabular}
\end{center}

\section{Automatic suggestions}

When the name of a variable or function resembles something that exists,
\hayashi{} suggests a correction:

\begin{lstlisting}[style=inline]
error: undefined function 'regres' — did you mean 'regress'?
\end{lstlisting}

Suggestions use edit distance (Levenshtein) over all names defined in scope
and in built-in functions.

\section{Re-throwing errors}

To propagate a caught error:

\begin{lstlisting}
try {
  risky_operation()
} catch e {
  display f"log: {e}"
  error(e)    // re-throws the error
}
\end{lstlisting}

The function \lstinline{error(msg)} throws a runtime error with the given
message.
